% Frugal AI seminar, part 2: why it matters for India. Figures come from the sources in the Ref lines;
% opinions are labelled as such. CheatSheet-safe: no floats, wide tables wrapped in adjustbox.

\section[India]{Why It Matters for India}

%%%%%%%%%%%%%%%%%%%%%%%%%%%%%%%%%%%%%%%%%%%%%%%%%%%
\begin{frame}[fragile]\frametitle{}
\begin{center}
{\Large Why It Matters for India}

\medskip
{\small (Languages, reach, data, and the cases that already work)}
\end{center}
\end{frame}

%%%%%%%%%%%%%%%%%%%%%%%%%%%%%%%%%%%%%%%%%%%%%%%%%%%
\begin{frame}[fragile]\frametitle{22 Languages, One Model?}
The Constitution's Eighth Schedule lists 22 languages:
\begin{center}
\footnotesize
\adjustbox{max width=\linewidth}{%
\begin{tabular}{llll}
Assamese & Bengali & Bodo & Dogri \\
Gujarati & Hindi & Kannada & Kashmiri \\
Konkani & Maithili & Malayalam & Manipuri (Meitei) \\
Marathi & Nepali & Odia & Punjabi \\
Sanskrit & Santali & Sindhi & Tamil \\
Telugu & Urdu & & \\
\end{tabular}}
\end{center}
\begin{itemize}
\item One general model rarely serves all of them equally well. The Hub report notes that a frontier generalist's English-centric training gives it ``no natural advantage'' on Indian-language tasks.
\item Smaller models built for one language or one job are a natural fit.
\end{itemize}
{\tiny (Ref: Eighth Schedule to the Constitution of India; Frugal AI Hub, Minds, Models \& Markets, 2026)}
\end{frame}

%%%%%%%%%%%%%%%%%%%%%%%%%%%%%%%%%%%%%%%%%%%%%%%%%%%
\begin{frame}[fragile]\frametitle{Languages Cost More Tokens}
\begin{itemize}
\item Sarvam found that Indian languages need more tokens than English for the same text.
\item Its co-founder is quoted: ``The same question, when asked in English, costs one-fifth of what it costs in an Indian language.''
\item Hosted models bill per token, and context memory and delay also grow with tokens, so more tokens means more money, more memory and more waiting.
\item A model that runs on your own computer removes the per-token bill.
\end{itemize}
{\tiny (Ref: Rest of World, Prabhu, Narayan, Sud, 7 April 2026. The figure is Sarvam's own.)}
\end{frame}

%%%%%%%%%%%%%%%%%%%%%%%%%%%%%%%%%%%%%%%%%%%%%%%%%%%
\begin{frame}[fragile]\frametitle{Reach: Phones, Voice, Low Bandwidth}
\begin{itemize}
\item Sarvam describes its aim as bringing generative AI to 800 million Indians with smartphones (the authors' figure).
\item Its work includes voice-enabled healthcare agents that run over WhatsApp and low-bandwidth connections.
\item AI4Bharat, launched at IIT Madras, builds lightweight systems that run on low-end phones and low-bandwidth networks.
\item Where the connection is slow or absent, a model on the device still works.
\end{itemize}
{\tiny (Ref: Rest of World, 7 April 2026)}
\end{frame}

%%%%%%%%%%%%%%%%%%%%%%%%%%%%%%%%%%%%%%%%%%%%%%%%%%%
\begin{frame}[fragile]\frametitle{Case: Soliga Voices}
\begin{itemize}
\item The Saving Voices Project, with IIIT Dharwad, built a text-to-speech system for the Soliga community in southern India, whose elders feared losing their language as younger members moved to cities.
\item It used five hours of voice data and runs offline for long periods on low-powered devices; the article names Raspberry Pi hardware running Linux, and the project targets ``sub-\$50 hardware''.
\item The voice data never left community devices, and community ownership was prioritised.
\item The trade-off, honestly stated: the word error rate is ``slightly high''.
\end{itemize}
{\tiny (Ref: Rest of World, R. Chandran, 2 April 2026)}
\end{frame}

%%%%%%%%%%%%%%%%%%%%%%%%%%%%%%%%%%%%%%%%%%%%%%%%%%%
\begin{frame}[fragile]\frametitle{Case: A Helper for Health Workers}
\begin{itemize}
\item An ASHA is, in the National Health Mission's words, ``a trained female community health activist'' chosen from the village, an ``interface between the community and the public health system'' and ``the first port of call for any health related demands''.
\item They counsel, help people reach services, and refer: a job of knowing the guidance and saying it clearly.
\item A frugal helper would sit on the worker's own phone, work offline, speak her language, look things up in reviewed guidance and refer. It would support her, never diagnose.
\item Status in the FrugalAI repository: set aside for now. The safety rules are written; nothing is built or tested for this case.
\end{itemize}
{\tiny (Ref: National Health Mission, ASHA page)}
\end{frame}

%%%%%%%%%%%%%%%%%%%%%%%%%%%%%%%%%%%%%%%%%%%%%%%%%%%
\begin{frame}[fragile]\frametitle{Case: Heritage on a Laptop, Modi Script}
\begin{itemize}
\item Modi is a historical script used for Marathi records. MoScNet reads handwritten Modi and writes Devanagari directly.
\item Its authors built a dataset of 2,043 Modi images with Devanagari transliterations.
\item By distillation, a student model about 163 times smaller (in parameters) than its teacher outperforms it, on their dataset. This is the authors' claim.
\item Why it fits frugal AI: old land and administrative records can be read on an ordinary computer, offline.
\end{itemize}
{\tiny (Ref: Kausadikar et al., Historic Scripts to Modern Vision, arXiv 2503.13060)}
\end{frame}

%%%%%%%%%%%%%%%%%%%%%%%%%%%%%%%%%%%%%%%%%%%%%%%%%%%
\begin{frame}[fragile]\frametitle{Indian Models and Efforts}
\begin{itemize}
\item Sarvam's 30-billion and 105-billion parameter models compete on Indian-language tasks rather than on the frontier's absolute scale (Hub report). Krutrim covers 22 Indian languages (Rest of World).
\item The Union minister for electronics and IT said in January 2026 that India is building 12 models of 50 to 120 billion parameters. That is his statement.
\item These are not small models. They are sovereign models, and small models complement them: the Principal Scientific Adviser's white paper calls small language models ``complementary components within a broader AI ecosystem'' (as reported by Storyboard18).
\end{itemize}
{\tiny (Ref: Frugal AI Hub 2026; Rest of World 2026; All India Radio, 21 January 2026; Storyboard18, 16 March 2026)}
\end{frame}
